\documentclass[11pt]{article}
\usepackage[margin=1in]{geometry}
\usepackage{amsmath,amssymb,amsthm}
\usepackage[T1]{fontenc}
\usepackage{lmodern}
\usepackage[hidelinks]{hyperref}
\newtheorem{theorem}{Theorem}
\title{Paths contradict the proposed acquaintance bound\\Conjecture 00000001679}
\author{gaochengzhecpu}
\date{October 9, 2026}
\begin{document}
\maketitle
\begin{abstract}
A path on $N$ vertices exposes only $N-1$ unordered pairs of agents at
each time. Even granting arbitrary rearrangements between times, complete
acquaintance requires at least $N/2-1$ rounds. For sufficiently large $N$
this exceeds $1024N/\log N$, disproving the stated universal upper bound.
\end{abstract}

\section{The actual acquaintance process}
Place one labeled agent at each vertex of the path $P_N$. Agents become
acquainted when their occupied vertices are adjacent. At each round the
usual game permits swaps along a matching of graph edges. Count adjacency
in the initial configuration as time zero, so $r$ rounds provide $r+1$
configurations. This convention gives the proposed upper bound its most
favorable starting point.

For the lower bound we allow an arbitrary permutation of the agents at
every time. This contains every legal matching-swap schedule and can only
make acquaintance easier. No assumption about which permutations can be
reached in a round is needed.

\begin{theorem}
If all $N$ agents on $P_N$ are acquainted after $r$ rounds, where $N\geq2$,
then
\[
 N\leq2(r+1),\qquad r\geq N/2-1.
\]
\end{theorem}
\begin{proof}
The path has $N-1$ edges. At each time these expose exactly $N-1$ unordered
pairs, or $2(N-1)$ ordered pairs of distinct agents. Repeated meetings add
no new pairs. The union of the ordered pairs exposed at times $0,\ldots,r$
therefore has size at most $2(r+1)(N-1)$. Complete acquaintance requires all
$N(N-1)$ ordered pairs of distinct agents. Hence
\[
 N(N-1)\leq2(r+1)(N-1).
\]
Canceling the positive factor $N-1$ proves the result.
\end{proof}

\section{Contradiction to the conjectured constant}
Choose an integer $N\geq4$ with $\log N>4096$; for example, any integer
larger than $e^{4096}+5$ will do. Then
\[
 \frac{1024N}{\log N}<\frac N4\leq\frac N2-1.
\]
Thus no schedule completing acquaintance on this connected path can have
$r\leq1024N/\log N$. This contradicts the conjectured bound
$a(P_N)\leq2^{10}N/\log N$. The proof is existential and exact, without
approximating $e^{4096}$ or evaluating a graph of that size.
The logarithm is natural, as usual; any other fixed base leads to the
same asymptotic contradiction.

\newpage
\section{Formal verification}
Lean represents the graph by Mathlib's actual
\texttt{SimpleGraph.pathGraph (n + 1)}. It proves that these paths are
connected. The formal parameter $n$ counts edges, so the paper's vertex
count is $N=n+1$. A schedule is a function from
\texttt{Fin (r + 1)} to permutations of \texttt{Fin (n + 1)}, assigning each
agent its actual position.

The proposition \texttt{AllMet} says that each two distinct agents occupy
adjacent vertices at some observed time. It is expressed using the
path graph's adjacency relation, rather than an independently invented
counting quantity. A finite map \texttt{visiblePair} enumerates each
time, each path edge, and its two orientations, and then transports the
edge's endpoints back to agents with the inverse position permutation.

\begingroup\raggedright
\begin{itemize}
\item \texttt{adjacent\_\allowbreak is\_\allowbreak visible} proves that every actual
meeting is in that finite image.
\item \texttt{counting\_\allowbreak lower\_\allowbreak bound} places all ordered
distinct-agent pairs inside the image, bounds its cardinality by the
domain's cardinality, and cancels $n>0$ to obtain $N\leq2(r+1)$.
\item \texttt{no\_\allowbreak short\_\allowbreak schedule} combines this inequality
with $N\geq4$, $\log N>4096$, and the conjectured upper bound to obtain a
contradiction.
\item \texttt{exists\_\allowbreak bad\_\allowbreak path} constructs a sufficiently
large natural vertex count through the Archimedean property and uses
the exact monotonicity of the real logarithm. It quantifies over every
round count satisfying the claimed bound and every position schedule.
\end{itemize}
\endgroup

The formal theorem is stronger than needed: even schedules with
unrestricted rearrangements fail within the claimed number of rounds.
Every actual adjacent-swap schedule supplies such a sequence of position
permutations. Accordingly the proof does not need to compute the exact
optimal acquaintance time or formalize an optimization algorithm.

\paragraph{Scope.}
The source gives a universal bound in terms of the number of vertices;
it imposes no expansion, density, or random-graph hypothesis. The examples
are ordinary finite connected paths. This disproof addresses that universal
upper bound; it makes no assertion about the best bounds for separately
restricted graph families.

\paragraph{Reproduction.}
In \texttt{lean/}, run:
\begin{verbatim}
lake build
lake env lean -DwarningAsError=true Main.lean
\end{verbatim}
The project pins Lean 4.19.0 and an exact Mathlib revision. The final
theorems use only standard logical axioms, with no unproved placeholders,
custom axioms, native evaluation, or external numerical certificates.
Run \texttt{tectonic main.tex} to rebuild the PDF. The bilingual source
and its immutable provenance are included with the submission.

\begingroup\raggedright
\begin{thebibliography}{9}
\bibitem{source} The Last Math Competition,
\href{https://github.com/The-Last-Math-Competition/The-Last-Math-Competition/blob/9b795e7a94a6076e49a65a6489e1caf9153abd23/conjectures/00000001679.md}{Conjecture 00000001679}, revision 9b795e7a94a6.
\bibitem{mathlib} The Mathlib Community,
\href{https://github.com/leanprover-community/mathlib4/commit/c44e0c8ee63ca166450922a373c7409c5d26b00b}{Mathlib, revision c44e0c8ee63c},
module \texttt{Combinatorics.SimpleGraph.Hasse}.
\end{thebibliography}
\endgroup
\end{document}
